\centerline{\bf \Huge Вектора}

\begin{flushright}
        \scriptsize{по ЕГЭ так по ЕГЭ...}
\end{flushright}

\problem{1}
Упростите выражение $\overline{C E}-\overline{C A}+\overline{E B}-\overline{D B}$.

\problem{2} \textbf{Передокажите задачу с занятия и навсегда запомните!}

(a) Пусть $A A_1$ - медиана треугольника $A B C$. Докажите, что

$$
\overline{A A_1}=\frac{1}{2}(\overline{A B}+\overline{A C})
$$

(b) На стороне $B C$ треугольника $A B C$ выбрана точка $A_1$ такая, что

$$
\frac{B A_1}{B C}=t
$$


Докажите, что

$$
\overline{A A_1}=(1-t) \overline{A B}+t \overline{A C}
$$

\problem{3}
В четырёхугольнике $A B C D$ точка $M$ - середина $B C, N$ - середина $A D$. Докажите, что

$$
M N \leqslant \frac{B A+C D}{2}
$$


Когда достигается равенство?

\problem{4}
Пусть $A B C D$ и $A B_1 C_1 D_1$ - два параллелограмма с общей вершиной. Докажите, что один из векторов $\overline{B B}_1, \overline{C C}_1$ и $\overline{D D}_1$ коллинеарен сумме двух других.

\problem{5}
Дан четырёхугольник $A B C D.\  A^{\prime}, B^{\prime}, C^{\prime}$ и $D^{\prime}$ - середины сторон $B C, C D, D A$ и $A B$ соответственно. Известно, что $A A^{\prime}=C C^{\prime}$ и $B B^{\prime}=D D^{\prime}$. Докажите, что $A B C D$ - параллелограмм.

\problem{6}
\textbf{[OMMO, 2016]}. В треугольнике $A B C$ с отношением сторон $A B: A C=5: 4$ биссектриса угла $B A C$ пересекает сторону $B C$ в точке $L$. Найдите длину отрезка $A L$, если длина вектора $4 \cdot \overrightarrow{A B}+5 \cdot \overrightarrow{A C}$ равна 2016 .